\chapter{Downloading, Auditing, and Importing Raw OULAD Data}
\label{ch:raw-import}

\begin{learningobjectives}
After completing this chapter, you will be able to:
\begin{itemize}
  \item explain the official archive, checksum, and expected-file contract;
  \item locate the extracted CSV files without committing raw data;
  \item explain why most raw columns initially land as text;
  \item trace the exact table creation and client-side bulk-load order;
  \item interpret import failure behavior and expected row counts; and
  \item reconcile the populated raw schema with saved source-audit evidence.
\end{itemize}
\end{learningobjectives}

\section{The raw-data contract}

The source data are downloaded on demand. They are not tracked in Git. The
expected local paths are:

\begin{lstlisting}
data/raw/oulad.zip
data/raw/source/assessments.csv
data/raw/source/courses.csv
data/raw/source/studentAssessment.csv
data/raw/source/studentInfo.csv
data/raw/source/studentRegistration.csv
data/raw/source/studentVle.csv
data/raw/source/vle.csv
data/raw/source/OULAD.names
\end{lstlisting}

The downloader in \projectfile{src/learning_analytics/download.py} verifies an
expected SHA-256 checksum and confirms that all eight archive members exist
before extraction.

\begin{projectconnection}
The exact archive checksum used by the implementation is stored as a constant
in \projectfile{src/learning_analytics/download.py}. The source auditor then calculates a checksum for
each extracted CSV and saves it in
\projectfile{reports/tables/source_file_profile.csv}.
\end{projectconnection}

\section{Download and source audit}

With the Python environment active:

\begin{lstlisting}[language=bash]
make data
\end{lstlisting}

This target runs two package commands:
\begin{enumerate}
  \item \texttt{python -m learning\_analytics download};
  \item \texttt{python -m learning\_analytics audit}.
\end{enumerate}

The source audit reads the large CSVs without loading every file into one
DataFrame. Its streaming profiler counts rows, records \texttt{?} or empty
cells, and calculates per-file checksums. Selected smaller files are then read
with pandas to calculate unique students, attempt keys, module counts, and
observed anomalies.

\begin{validationbox}
After \texttt{make data}, require:
\begin{lstlisting}[language=bash]
test -f data/raw/source/studentVle.csv
test -f reports/source-audit.json
test -f reports/tables/source_file_profile.csv
\end{lstlisting}
Then compare the profile with \cref{tab:source-files}. Do not proceed if file
names, checksums, or counts differ unexpectedly.
\end{validationbox}

\section{Why raw tables preserve source shape}

The raw schema is a landing zone. In
\projectfile{sql/01_raw_tables.sql}, most columns are declared as
\texttt{text}. This design has three advantages:

\begin{enumerate}
  \item values such as \texttt{?} can load before null conversion;
  \item source contents remain available for reconciliation; and
  \item type errors are handled visibly in the staging layer.
\end{enumerate}

The largest table, \texttt{raw.student\_vle}, uses integer columns already and
is declared \texttt{UNLOGGED} for load and analytical performance. Unlogged
tables reduce write-ahead logging but are not designed for durable source
storage across server crashes. The official CSV remains the reproducible
source.

\begin{lstlisting}[language=SQL,caption={Selected raw landing definitions from
\texttt{sql/01\_raw\_tables.sql}.}]
CREATE TABLE raw.assessments (
    code_module text,
    code_presentation text,
    id_assessment text,
    assessment_type text,
    date text,
    weight text
);

CREATE UNLOGGED TABLE raw.student_vle (
    code_module text,
    code_presentation text,
    id_student integer,
    id_site integer,
    date integer,
    sum_click integer
);
\end{lstlisting}

\section{Import order}

Run:

\begin{lstlisting}[language=bash]
make ingest
\end{lstlisting}

The target delegates to \projectfile{scripts/load_raw.sh}. The script:
\begin{enumerate}
  \item creates the six schemas;
  \item drops and recreates the seven raw tables;
  \item loads each CSV with PostgreSQL \texttt{\textbackslash copy};
  \item creates two supporting raw indexes; and
  \item runs \texttt{ANALYZE} on the largest raw relations.
\end{enumerate}

\subsection{Why \texttt{\textbackslash copy} is used}

\texttt{\textbackslash copy} is a \texttt{psql} client command. It reads a file
visible to the client process. Server-side \texttt{COPY} reads a path visible
to the PostgreSQL server and may require different permissions. Because the
script runs from the project machine and constructs absolute project paths,
client-side loading is appropriate.

\begin{lstlisting}[language=SQL,caption={The project bulk-load pattern.}]
\copy raw.courses
FROM '/absolute/project/path/data/raw/source/courses.csv'
WITH (FORMAT csv, HEADER true)
\end{lstlisting}

The course shows a descriptive absolute path above, but the implementation
calculates the real path at runtime. Do not hard-code a personal home directory
in the project or course source.

\section{Failure behavior}

The shell script begins with:

\begin{lstlisting}[language=bash]
set -euo pipefail
\end{lstlisting}

Every \texttt{psql} call uses:

\begin{lstlisting}[language=bash]
-v ON_ERROR_STOP=1
\end{lstlisting}

Together, these settings stop the load when a command fails rather than
continuing into an apparently complete database.

\begin{editionnote}
\textbf{Suggested enhancement, not implemented:} the project does not provide a
rejected-row quarantine table. A malformed CSV row stops the bulk load.
Edition 0.1 teaches how to inspect the error and source file; a future
production-oriented module could add staged quarantine handling.
\end{editionnote}

\section{Validate the imported row counts}

\begin{lstlisting}[language=SQL,caption={Instructional row-count reconciliation
query based on the implemented gallery.}]
SELECT 'courses' AS source_table, count(*) AS rows
FROM raw.courses
UNION ALL
SELECT 'assessments', count(*) FROM raw.assessments
UNION ALL
SELECT 'vle', count(*) FROM raw.vle
UNION ALL
SELECT 'student_info', count(*) FROM raw.student_info
UNION ALL
SELECT 'student_registration', count(*)
FROM raw.student_registration
UNION ALL
SELECT 'student_assessment', count(*)
FROM raw.student_assessment
UNION ALL
SELECT 'student_vle', count(*) FROM raw.student_vle
ORDER BY source_table;
\end{lstlisting}

The output grain is one row per source table. The sum should be 10,900,970.
The expected individual counts are listed in \cref{tab:source-files}.

\begin{validationbox}
A correct import requires more than ``seven tables exist.'' Check:
\begin{enumerate}
  \item all expected tables exist;
  \item each row count matches the source audit;
  \item no source file was loaded into the wrong table;
  \item the source profile and database counts reconcile; and
  \item the next staging step can cast expected values.
\end{enumerate}
\end{validationbox}

\section{From raw text to staging types}

The next pipeline stage creates typed views. For example:

\begin{lstlisting}[language=SQL]
SELECT
    trim(code_module) AS code_module,
    id_assessment::integer AS id_assessment,
    NULLIF(date, '?')::integer AS due_day,
    weight::numeric(6, 2) AS weight
FROM raw.assessments;
\end{lstlisting}

Read the expressions from inside out:
\begin{enumerate}
  \item \texttt{NULLIF(date, '?')} converts the missing marker to SQL null;
  \item \texttt{::integer} casts remaining values to integers;
  \item aliases replace source names with analytical names; and
  \item \texttt{numeric(6,2)} preserves assessment weights with decimal scale.
\end{enumerate}

If an unexpected nonnumeric value appears, the cast fails. That is useful:
the source contract has changed or the input is wrong.

\section{Knowledge check}

\begin{enumerate}
  \item What does the archive checksum verify?
  \item Why are most raw columns text?
  \item What is the difference between \texttt{COPY} and
    \texttt{\textbackslash copy}?
  \item What does \texttt{ON\_ERROR\_STOP} protect against?
  \item Why is table existence insufficient import validation?
  \item What tradeoff comes with an unlogged table?
\end{enumerate}

\begin{exercisebox}{Exercise 6.1: Reconcile one imported table}
Choose \texttt{studentAssessment.csv}. Produce a short reconciliation containing:
\begin{enumerate}
  \item source path;
  \item source-profile row and column counts;
  \item raw database row count;
  \item raw column names;
  \item candidate key;
  \item known missing marker rule; and
  \item one query that would reveal duplicate candidate keys.
\end{enumerate}
\solutionref{sol:reconcile}
\end{exercisebox}

\begin{exercisebox}{Exercise 6.2: Explain a failed cast}
Suppose one \texttt{assessments.date} value is the text \texttt{TBD}. Explain:
\begin{enumerate}
  \item why the raw load can succeed;
  \item why the staging cast fails;
  \item how to locate the source row;
  \item why silently converting every invalid value to null is risky; and
  \item what evidence should be recorded before deciding on a rule.
\end{enumerate}
\solutionref{sol:failed-cast}
\end{exercisebox}

\begin{commonmistakes}
\begin{itemize}
  \item Running import before download and extraction.
  \item Loading the same CSV twice without recreating raw tables.
  \item Confusing \texttt{studentInfo} attempts with unique students.
  \item Treating a successful bulk load as proof of correct column mapping.
  \item Hard-coding a personal absolute path in SQL.
  \item Changing a source value to make a cast pass without documenting the
    evidence.
\end{itemize}
\end{commonmistakes}

\transferheading
The same raw-import pattern works for many datasets: verify the acquired file,
preserve a source-faithful landing layer, stop on structural failures, cast in
a visible staging layer, and reconcile row counts before analysis.

\begin{summarybox}
The project downloads and checksum-verifies the official OULAD archive, profiles
the extracted sources, creates source-faithful raw tables, and bulk-loads seven
CSV files with fail-fast client-side commands. Correctness requires
reconciliation of files, columns, rows, types, and missing markers, not merely a
successful command.
\end{summarybox}
